\medskip

Dans un laboratoire, on étudie une réaction chimique dans un réacteur fermé, sous certaines conditions. Le traitement numérique des données expérimentales a permis de modéliser l'évolution de la température de cette réaction chimique en fonction du temps.

L'objectif de cet exercice est d'étudier cette modélisation.

La température est exprimée en degré Celsius et le temps est exprimé en minute.

Dans tout l'exercice, on se place sur l'intervalle de temps $[0~;~10]$.

Les parties A et B peuvent être traitées de façon indépendante.

\medskip

\textbf{Partie A}

\medskip

Dans un repère orthogonal du plan, on donne ci-dessous la courbe représentative de la fonction température en fonction du temps sur l'intervalle $[0~;~10]$.

\medskip

\begin{enumerate}
\item Déterminer, par lecture graphique, au bout de combien de temps la température redescend à sa valeur initiale à l'instant $t = 0$.
\end{enumerate}

\begin{center}
\psset{xunit=1cm,yunit=0.1cm,arrowsize=2pt 3}
\begin{pspicture}(-1,-5)(10,73)
\multido{\n=0+1}{11}{\psline[linewidth=0.15pt](\n,0)(\n,70)}
\multido{\n=0+10}{8}{\psline[linewidth=0.15pt](0,\n,)(10,\n)}
\psaxes[linewidth=1.25pt,labelFontSize=\scriptstyle,Dy=10]{->}(0,0)(0,0)(10,70)
\psaxes[linewidth=1.25pt,labelFontSize=\scriptstyle,Dy=10](0,0)(0,0)(10,70)
\psplot[plotpoints=2000,linewidth=1.25pt,linecolor=red]{0}{10}{60 x mul 40 add 2.71828 0.5 x mul exp div}
\uput[u](9.5,0){$t$ (min)}\uput[r](0,71){$\degres$ C}
\end{pspicture}
\end{center}

On appelle $f$ la fonction température représentée par la courbe ci-dessus.

On précise que la fonction $f$ est définie et dérivable sur l'intervalle $[0~;~10]$.

On admet que la fonction $f$ peut s'écrire sous la forme $f(t) = (at + b)\e^{-0,5 t}$ où $a$ et $b$ sont deux constantes réelles.

\begin{enumerate}[start=2]
	\item On admet que la valeur exacte de $f(0)$ est 40. En déduire la valeur de $b$.
	\item On admet que $f$ vérifie l'équation différentielle $(\mathrm{E}): y'+ 0,5 y= 60 \e^{-0,5 t}$. Déterminer la valeur de $a$.
\end{enumerate}
\medskip 

\textbf{Partie B : Étude de la fonction $f$}
\medskip 

On admet que la fonction $f$ est définie pour tout réel $t$ de l'intervalle $[0~;~10]$ par

\[f(t) = (60 t+ 40)\e^{-0,5 t}\]

\smallskip

\begin{enumerate}
\item Montrer que pour tout réel $t$ de l'intervalle $[0 ; 10]$, on a : $f'(t)=(40-30 t) \e^{-0,5 t}$.
\item
	\begin{enumerate}
		\item Étudier le sens de variation de la fonction $f$ sur l'intervalle $[0 ; 10]$.


Dresser le tableau de variations de la fonction $f$ en y faisant figurer les images des valeurs présentes dans le tableau.

		\item Montrer que l'équation $f(t) = 40$ admet une unique solution $\alpha$ strictement positive sur l'intervalle $]0~;~10]$.
		\item Donner une valeur approchée de $\alpha$ au dixième près et en donner une interprétation dans le contexte de l'exercice.
	\end{enumerate}
\item On définit la température moyenne, exprimée en degré Celsius, de cette réaction chimique entre deux temps $t_{1}$ et $t_{2}$, exprimés en minute, par

\[\dfrac{1}{t_{2}-t_{1}} \int_{t_{1}}^{t_{2}} f(t)\,\rm{d} t\]

	\begin{enumerate}
		\item À l'aide d'une intégration par parties, montrer que

\[\int_{0}^{4} f(t) \, \rm{d}t=320-\dfrac{800}{\e^{2}}\]

		\item En déduire une valeur approchée, au degré Celsius près, de la température moyenne de cette réaction chimique au cours des 4 premières minutes.
	\end{enumerate}
\end{enumerate}
